% ============================================================================
\section{Settings and notation}
\label{app:settings}

This appendix lists the settings of the method and tests them one by one, then gathers the notation of the paper.
\Cref{tab:settings} lists the settings, none of which is tuned on the mission labels.

\begin{table}[!htbp]
\centering
\caption{Fixed settings of the method. A \emph{resolution} setting can be refined without changing
what is computed, only how finely; a \emph{design} setting defines what is computed.}
\label{tab:settings}
\tabin{settings}
\end{table}

To check that no setting decides the result, each one was varied alone on Top-1, through the whole method. Each
version was compared without labels with a reference: its finest value for a resolution setting, and the value of
the method for a design choice. The
comparison uses the rank correlation of the distances, field by field and scale by scale, and the agreement
of the taxonomies (\cref{tab:study}). None of them changed the number of behaviors, and at most
\nm{study-runs-moved} of the $250$ runs changed cluster (an adjusted Rand index of at least
\nm{study-ari-min} with the reference). The distances depend mostly on how finely the domain is sampled at
the smallest scale, which is why the seeds follow the scale (\cref{sec:seeds}).

\begin{table}[!htbp]
\centering
\caption{Each setting varied alone on Top-1, without labels. \emph{Distances}: the Spearman rank correlation
between the distances at the value of the method and at the finest value tested, the smallest over
the twenty matrices of fields and scales. For a design choice, it is the smallest over the values tested.
\emph{Taxonomy}: the smallest adjusted Rand index between the consensus partition of a value tested and
that of the reference. Resolution settings are checked on the domain of the method (seeds within
$\sigma/2$ of every robot position) when it was run, otherwise on a fixed subset of $\nm{n-coreset}$ robot positions.}
\label{tab:study}
\tabin{study}
\end{table}

The notation of the paper is gathered in \cref{tab:notation}, grouped by the step of the method that introduces it.

\begin{table}[!htbp]
\centering
\caption{Notation, grouped by the step of the method that introduces it.}
\label{tab:notation}
\begin{tabular*}{\linewidth}{@{\extracolsep{\fill}}l>{\raggedright\arraybackslash}p{6.4cm}l@{}}
\toprule
Symbol & Meaning & Defined in\\
\midrule
\multicolumn{3}{@{}l}{\emph{Recordings and fields}}\\
$N$, $T$, $\Delta t$ & agents, instants per recording, sampling step & \cref{sec:fields}\\
$\vx_i(t)$, $\vv_i(t)$ & position and forward-difference velocity of agent $i$ (vectors) & \cref{sec:fields}\\
$O_i$, $\sigma$ & Gaussian footprint and its width, the observation scale & \eqref{eq:footprint}\\
$M_0$, $\vM_1$, $\tM_2$ & moments of order zero, one and two of the velocity distribution (scalar, vector, tensor) & \eqref{eq:moments}\\
$\vV_\sigma$ & flux, the vector field $\vM_1$ & \eqref{eq:hhd}\\
$\rho_\sigma$, $K_\sigma$ & density and kinetic intensity & \cref{sec:fields}\\
$\varphi_\sigma$, $\psi_\sigma$ & divergence and vorticity potentials of the Helmholtz--Hodge decomposition & \eqref{eq:potentials}\\
$\kappa_\sigma$ & radial kernel of the potentials & \eqref{eq:kappa}\\
$\vbar$, $\overline{v^2}$ & mean speed and mean squared speed of the recording & \cref{sec:fields}\\
\addlinespace
\multicolumn{3}{@{}l}{\emph{Persistence and distances}}\\
$\sigma_1,\dots,\sigma_M$ & ladder of scales, $M=4$ & \eqref{eq:ladder}\\
$J_\sigma$ & Jacobi coupling $\sigma^2\det(\nabla\varphi_\sigma,\nabla\psi_\sigma)$ & \eqref{eq:jacobi}\\
$D^{s}_{f,m}(A)$ & diagram $s$, one of six, of field $f$ at scale $\sigma_m$ for recording $A$ & \eqref{eq:sw}\\
$d_{f,m}$ & sliced Wasserstein distance between recordings & \eqref{eq:sw}\\
\addlinespace
\multicolumn{3}{@{}l}{\emph{Fusion and clustering}}\\
$\tilde d_{f,m}$ & percentile rank of the mutual rank of $d_{f,m}$ & \cref{sec:fusion}\\
$W_m$ & balanced view at scale $\sigma_m$ & \eqref{eq:view}\\
$\knn$, $t$ & neighborhood size and diffusion iterations of the fusion & \eqref{eq:snf}\\
$\Sstar$ & consensus network after fusion & \cref{sec:fusion}\\
$C$, $\Cstar$ & number of clusters, its modal value across the local grid & \cref{sec:ap}\\
\bottomrule
\end{tabular*}
\end{table}

% ============================================================================
\section{The ladder of scales}
\label{app:gallery}

\Cref{fig:scales_all} shows the five fields of the Shelter exemplar of Top-1 at the four scales of
the ladder. At the finest scale the fields resolve individual paths and the walls of the shelter; at
the coarsest they keep the organization of the whole group.

\begin{figure}[!htbp]
\centering
\figin{0.70\linewidth}{figures/fig_scales.pdf}
\caption{The five fields of one recording (the Shelter exemplar of Top-1) at the four scales of the
ladder, $\nm{sigma1-cm}$, $\nm{sigma2-cm}$, $\nm{sigma3-cm}$ and $\nm{sigma4-cm}$\,cm.
Color scales are per panel: sequential for the density $\rho_\sigma$ and the kinetic intensity
$K_\sigma$ (zero is pale blue), diverging for the potentials $\varphi_\sigma$, $\psi_\sigma$ and the
coupling $J_\sigma$ (blue negative, orange positive, zero gray). White marks points farther than
$5\sigma$ from every robot at every instant, where the fields are not evaluated.}
\label{fig:scales_all}
\end{figure}
\FloatBarrier

The top of the ladder follows from the merging width $D_s/2.35$ at which the footprints of the whole
group stop being resolved (\cref{sec:ladder}). \Cref{fig:uncapped} shows what lies near it. On a ladder
extended up to that width, the density and the vorticity potential of the Homing exemplar keep their
organization up to a few tens of centimeters. Past that, each field becomes one smooth bump over the
whole group, and its diagrams no longer depend on how the robots moved. The cap of the method therefore lies far
below the merging width. $D_s$ measures the extent of the swarm, not the size of the structures of its
motion, such as the corridor, the shelter or a group of robots.

\begin{figure}[!htbp]
\centering
\figin{\linewidth}{figures/figA_uncapped.pdf}
\caption{The fields near the merging width. The density and the vorticity potential of the Homing
exemplar on a ladder that extends to the merging width of the whole group, $D_s/2.35$. Beyond a few
tens of centimeters every field becomes a single featureless bump, whose diagrams no longer
distinguish behaviors.}
\label{fig:uncapped}
\end{figure}
\FloatBarrier

The method uses the cap $k=\nm{cap-k}$, and its taxonomy hardly depends on this value. For caps
$\sigma_{\max}=D_s/(2.35\,k)$ from a third of the merging width down to a twentieth of it, the whole method was run on
Top-1 and the taxonomy obtained as in \cref{sec:ap}. Every cap from $k=\nm{cap-same-kmin}$ to $k=\nm{cap-same-kmax}$, a
top of the ladder between $\nm{cap-same-cmmin}$ and $\nm{cap-same-cmmax}$\,cm, gives the same taxonomy, and
$k=\nm{cap-apart-k}$ moves \nm{cap-apart-runs} run. Even with no cap at all, on a fixed subset of $\nm{n-coreset}$ robot
positions, the taxonomy differs from that of the capped ladder by \nm{uncapped-runs-moved} run, and adding scales between
the same bounds, up to $\nm{uncapped-mmax}$, leaves it unchanged.

% ============================================================================
\section{Robustness of the fusion}
\label{app:fusion}

\subsection{Alternative fusion architectures}
\Cref{tab:designs} compares, on Top-1 and over the same $99$ settings, our fusion with
alternatives built from the same rank-normalized matrices and clustered by the same affinity
propagation. Without diffusion, affinity propagation splits the runs of each single controller of
Top-1 into several groups, $\nm{design-none-C}$ in all, although each mission is one controller there,
and every setting returns the same partition. Diffusing first among the
fields of each moment, then across scales, then across moments, compounds the smoothing. The pair $(\knn,t)$
no longer means a single resolution, missions merge already at moderate settings (median purity
\nm{design-hierarchical-medpur}), and the number of clusters stops decreasing regularly with $t$.
A single SNF
over all field-and-scale views gives the flux three times the weight of the density. Among the
architectures that diffuse, the two fusions that balance the moments and the single SNF over twelve views, one per
moment and scale, reach a median purity of at least \nm{design-top-medpur-min}, and the two others of at most
\nm{design-low-medpur-max}. The settings in
which purity and ARI are both one number \nm{design-balanced-perfect} for our fusion, \nm{design-fluxrerank-perfect}
for the fusion that re-ranks the flux first, and \nm{design-snf12-perfect} for the single SNF over twelve views.
Our fusion keeps the number of clusters from increasing with $t$ at \nm{design-balanced-mono} of the nine
neighborhood sizes, the re-ranked fusion at \nm{design-fluxrerank-mono}, and the single SNF over twelve views at
\nm{design-snf12-mono}.

\begin{table}[!htbp]
\centering
\caption{Fusion architectures on Top-1 (same distances, same clustering, $99$ settings).
\emph{Modal}: settings at the modal count $\Cstar$; \emph{perfect}: settings with purity and ARI equal to
one; $C\le2$: settings where the partition collapses; \emph{monotone}: neighborhood sizes for which
$C$ never increases with $t$.}
\label{tab:designs}
\tabin{designs}
\end{table}


\subsection{Rank fusion against standard SNF}
On Top-1, with $\knn=5$, rank fusion settles on \nm{std-ours-C} clusters from $t=\nm{std-ours-t-from}$ on and
keeps them up to $t=\nm{std-t-max}$. With the affinity of standard SNF on the raw distances instead of ranks, the
number of clusters drifts with the number of diffusion iterations, from \nm{std-C-first} at
$t=\nm{std-t-first}$ to \nm{std-merge-C} from $t=\nm{std-merge-from}$ on, where two missions share a cluster
(purity \nm{std-merge-purity}).

\subsection{Convergence of affinity propagation}
With damping $0.9$ the messages of affinity propagation change slowly, so the algorithm must be
given enough iterations to settle. The settings of the original implementation declare convergence
after a window of $100$ unchanged iterations, within at most $1000$. With them it converges at every
setting of Top-1 and of Top-10, the largest dataset, after a median of \nm{ap-top10-iter-median}
iterations on Top-10 and at most \nm{ap-top10-iter-max}. A window of $200$ or $400$ iterations, or a ceiling of $5000$, leaves
each of the $99$ partitions of both datasets unchanged: the algorithm has time to finish. The
shorter window of scikit-learn, $15$ iterations, stops earlier and changes \nm{ap-short-top1-changed}
partitions of Top-1 and \nm{ap-short-top10-changed} of Top-10. Another damping, $0.8$ or $0.95$, or
another draw of the tie-breaking noise can reach another set of exemplars, at up to
\nm{ap-other-changed-max} of the $99$ settings. The modal consensus taxonomy stays the same,
with an adjusted Rand index of at least \nm{ap-other-cons-min} with the reference.

\subsection{Each step of the fusion changed alone}
\begin{table}[!htbp]
\centering
\caption{Each step of the fusion changed alone, on the ten datasets, the taxonomy of every variant
built on its own modal region, as for the method. \emph{Modal} and \emph{stability}: means over the datasets, without
labels. \emph{Purity} and \emph{ARI}: means over the datasets. \emph{Misassigned}: runs outside a
cluster dominated by their mission, summed over the datasets.}
\label{tab:variants}
\tabin{variants}
\end{table}

\Cref{tab:variants} changes one step of the fusion at a time and keeps the others. On the ten
datasets, the taxonomy of every variant is built on its own modal region, as for the method. The first step is the affinity
given to each view. Percentile ranks alone, without the mutual ranks, misassign
\nm{fx-affpercentile-wrong} runs, \nm{fx-affpercentile-xs} of them XOR Aggregation runs placed with
Shelter (\cref{app:merge}). Percentile ranks passed through the normal quantile function misassign
\nm{fx-affprobit-wrong}. A Gaussian kernel whose width adapts to each pair, through the distances of the two runs to their seventh neighbors, keeps the purity. It splits the missions
further, however, with a mean adjusted Rand index of \nm{fx-affselftuning-ari}. Inside one scale, an equal weight per field misassigns
\nm{fx-fieldsequal-wrong} runs. The weights can also be read from the data. The spectral gap of a
field at one scale is the largest gap between two consecutive eigenvalues of the normalized Laplacian of its
neighbor graph, the first gap left out \citep{von2007tutorial}. This graph joins two runs when one is among the
seven nearest neighbors of the other, with the weight $e^{-d_{f,m}/s}$, where $s$ is the standard deviation of
the entries of $d_{f,m}$. A larger gap means that the field splits the runs into clearer groups. Given to each
field alone, these gaps misassign \nm{fx-fieldsconfidence-wrong} runs. Shared within each moment, that is,
divided by the number of fields of the moment and normalized to sum to one, they give the weights read from the
data that \cref{sec:fusion} mentions. They reduce to the weights of \eqref{eq:view} when the gaps are equal, and they misassign
\nm{fx-fieldssourceconf-wrong} runs. The
weighted geometric mean misassigns \nm{fx-fieldsgeomean-wrong} runs, and the maximum over the fields
\nm{fx-fieldsmax-wrong}. Across scales, the diffusion is needed: averaging the four views without it returns one partition at every setting, which splits the
missions into \nm{fx-mean-cmin} to \nm{fx-mean-cmax} groups as controllers are pooled. A single SNF over
the twenty field-and-scale views merges missions, with \nm{fx-scalessnf20-wrong} runs misassigned. On
the consensus, affinity propagation with the minimum preference reproduces the taxonomy of the method, with the
same $\Cstar$ on every dataset and an adjusted Rand index of at least \nm{fx-apmin-ari-min} with it. Louvain
modularity finds the same $\Cstar$ on \nm{fx-louvain-same} of the ten datasets, at an adjusted Rand index of at least
\nm{fx-louvain-ari-min}. Clusterings that choose the number of groups from an
eigengap, from the largest gap of an average-linkage tree or from density merge missions instead,
into \nm{fx-collapse-cmin} to \nm{fx-collapse-cmax} groups, with a purity of at most
\nm{fx-collapse-pmax}. A higher adjusted Rand
index in the table goes with fewer groups on the pooled datasets, not with fewer misassigned runs,
since that index counts the sub-behaviors against the method (\cref{sec:protocol}).


% ============================================================================
\section{Other descriptions on all the runs}
\label{app:baselines}

\Cref{tab:baselines_full} compares the descriptions of \cref{tab:baselines} on Top-1 and on Top-10, where the missions
split into sub-behaviors. There, homogeneity and completeness are the scores to read, since the ARI counts the
sub-behaviors as errors.

\begin{table}[!htbp]
\centering
\caption{Descriptions of the runs under the same fusion, affinity propagation and modal consensus, on Top-1 and Top-10. The
number of modal settings follows $\Cstar$ in parentheses; \emph{error-free}: settings, out of $99$, in which no run is
misassigned, those in the modal region in parentheses. The method is shaded. $^{*}$Partition published for the same Top-1
runs, with five fixed clusters \citep{legarda2026universal}.}
\label{tab:baselines_full}
{\footnotesize\renewcommand{\arraystretch}{1.0}\tabin{baselines}}
\end{table}
\FloatBarrier

% ============================================================================
\section{Rotations on the robot data}
\label{app:invariance}

To test the invariance on the robot data, five runs of Top-1, one per mission, were rotated by random angles about the center of the arena
and processed again with the same code. For every field and scale, the distance between each run and
its rotated copy was divided by the median distance between the run and twelve other runs of its mission. The
density, the kinetic intensity and both potentials are recovered up to rounding, within \nm{inv-other-max} of that
median distance. The coupling differs most at the finest scale, by at most \nm{inv-worst} of it, because it is the
most sensitive to the placement of the seeds.

% ============================================================================
\section{XOR Aggregation runs that resemble Shelter}
\label{app:merge}

\begin{figure}[!htbp]
\centering
\figin{\linewidth}{figures/figA_merge.pdf}
\caption{XOR Aggregation runs that percentile ranks alone, without mutual ranks, place in the Shelter cluster on
Top-4. (a)~Median fraction of robots
moving faster than $1$\,cm/s over the run, for these runs, for the other XOR runs and for the Shelter
runs. (b--d)~Final positions of the robots, pooled over the runs of each group.}
\label{fig:merge}
\end{figure}

With percentile ranks alone instead of mutual ranks, XOR Aggregation runs join the Shelter cluster. On Top-3, all
\nm{merge-top3-n} XOR Aggregation runs do, and on Top-4, \nm{merge-top4-n} of the $200$ do, from every controller of the
mission. The method places none of them with Shelter on either dataset. Compared with the other XOR runs of Top-4, these
runs settle late, after a median of \nm{merge-top4-settle-abs}\,s against \nm{merge-top4-settle-kept}\,s. They also move
about twice as fast on average (\nm{merge-top4-speed-abs} against \nm{merge-top4-speed-kept}\,m/s, and
\nm{merge-top4-speed-shel}\,m/s for Shelter), yet they end on the same two black areas (\cref{fig:merge}). What keeps them
with their mission is their neighborhood. Among their seven nearest neighbors in each view, Shelter runs make up
\nm{merge-top4-nn-pct}\,\% with percentile ranks, against \nm{merge-top4-nn-mut}\,\% with mutual ranks. The Shelter runs
have one another as nearest neighbors, so their mutual ranks with these runs are large.

% ============================================================================
\section{Stability and reproducibility of the attribution}
\label{app:attribution}

The attribution holds the seeds fixed, which is justified at the critical points inside the domain.
There the gradient of the field vanishes, so that a critical value does not change at first order when
its point moves with the trajectories, as the envelope theorem states. On a synthetic milling group,
perturbing the tracks changes the minimum of $\psi_\sigma$ by the same amount, to the precision of the
finite differences, whether the minimum is relocated or held in place. At a seed that
moves with a robot, as the seeds taken from the robot positions do, the two derivatives differ even in sign
(\nolinkurl{manuscript/verify_propositions.py}). The value distance splits in the same way as the sliced
distance: the term of each quantile is shared equally between the two runs, and the half of a run between the two
order statistics that the quantile interpolates.

The split of \cref{sec:res_attribution} was checked in two ways. First, a Gaussian noise was added to every
position of the Directional Gate exemplar, with the seeds held fixed. A noise of $2$\,mm changes the field at the
seeds by up to \nm{attr-stab-field-change} of its largest value. It leaves \nm{attr-stab-pers} of the topological
split and \nm{attr-stab-val} of the value split on the same seeds, and the topological split is the more stable at
every noise tested (\nm{attr-stab5-pers} against \nm{attr-stab5-val} at $5$\,mm). The term of a quantile falls on
the seed that holds that rank, which a small noise can change, while the share of a diagram point stays on its
critical points. Second, the split was repeated over the \nm{attr-runs} runs of the Directional Gate controller,
each against the same Foraging exemplar. For the persistence diagrams, the two extrema with the largest shares are
a maximum and a minimum of $\psi_\sigma$, the two rotation centers of the traffic, \nm{attr-centres-phrase}
\nm{attr-runs} runs, with a median share of \nm{attr-centres-share}. Half of the mean map lies in
\nm{attr-mean-half-pers} of the arena, against \nm{attr-mean-half-val} for the value distributions, whose split
covers the whole positive cell of the potential. The centers themselves move by about \nm{attr-centres-spread} from
run to run, so two single topological maps overlap less on average (\nm{attr-overlap-pers}) than two diffuse value
maps (\nm{attr-overlap-val}). What recurs from run to run is the structure, not its exact position.

The split for the ten pairs of missions is given in \cref{tab:attribution} (\cref{sec:res_attribution}). In
\nm{attrp-pairs-50} of the ten pairs, the two largest extrema are a maximum and a minimum in every one of the
\nm{attr-runs} runs; only Directional Gate against XOR Aggregation drops to $\nm{attrp-recur-min}/50$, because one run
splits its largest weights across two maxima of $\rho_\sigma$ before reaching its minimum.

